\chapter{イベント・ドリブン}
\label{ch:event}

企業の出来事（決算、アナリストの推奨の変更、合併、指数の入れ替え、破綻、アクティビストの介入）の前後で生じる価格の歪みをとる戦略である。LLM-BLM の研究に最も近いのが、決算発表の後のドリフトなので、そこから始める。

\section{決算発表の後のドリフト}

\subsection{発見}

Ball・Brown（1968）\cite{ballbrown1968} は、会計の利益の情報が株価に反映されることを示した最初の研究の一つで、すでに、利益が予想を上回った会社の株価が、発表の後もしばらく上がり続ける傾向を観察していた。

\begin{paperbox}[Bernard・Thomas（1989, 1990）\cite{bernardthomas1989,bernardthomas1990}]
四半期の利益の驚き（SUE：前年同期からの変化を、その変動の大きさで割ったもの）で銘柄を 10 の組に分けると、驚きの大きい組は発表の後 60 取引日にわたって上がり続け、小さい組は下がり続ける（決算発表の後のドリフト、PEAD）。この遅れた反応は、リスクの対価では説明しにくい。1990 年の論文は、四半期の利益の驚きには自己相関がある（今期が良ければ、来期も前年同期より良いことが多い）のに、株価はそれを十分に織り込まず、次の四半期の発表のときに、予測できる方向の反応が起きることを示した。投資家が、利益の時系列の性質を単純化して捉えている、という解釈である。
\end{paperbox}

\subsection{注目と反応の遅れ}

PEAD は、投資家の注目が足りないときに大きくなる。

\begin{itemize}
\item \textbf{金曜日の発表}：DellaVigna・Pollet（2009）\cite{dellavigna2009} は、金曜日に発表された決算では、発表の直後の反応が小さく、その後の遅れた反応が大きいことを示した。週末の前で、投資家の注目が下がるためと解釈される。
\item \textbf{同じ日の発表の多さ}：Hirshleifer・Lim・Teoh（2009）\cite{hlt2009} は、同じ日にほかの会社の決算発表が多いほど、直後の株価と売買高の反応が小さく、ドリフトが大きいことを示した（注意の分散）。競合するニュースの少ない日の発表では、PEAD を狙う戦略は利益にならなかった。
\item \textbf{情報の処理の遅れ}：Akbas・Markov・Subasi・Weisbrod \cite{akbas2017} は、アナリストの予想のデータベースへの反映が遅れる銘柄ほど、発表のときの反応が小さく、その後のドリフトが大きいことを示した。
\end{itemize}

\subsection{流動性と、消えたドリフト}

\begin{paperbox}[Chordia・Goyal・Sadka・Sadka・Shivakumar（2009）\cite{cgsss2009}]
PEAD は、主に流動性の極めて低い銘柄で起きていることを示した。驚きの大きい株を買い小さい株を売る戦略の時価総額加重のリターンは、最も流動性の高い銘柄で月 0.14\%、最も低い銘柄で月 1.60\% だった。流動性の低い銘柄は売買のコストとインパクトが大きく、さまざまなコストの推定で、紙の上の収益の 63〜100\% がコストで消える。コストそのものがドリフトの残る原因かもしれない、と論じた。
\end{paperbox}

\begin{paperbox}[Martineau（2022）「安らかに眠れ、決算発表の後のドリフト」\cite{martineau2022}]
現代の米国株では、決算の驚きは発表の当日に株価に完全に反映され、PEAD は消えた、と報告した。大型株では 2006 年以降 PEAD は見られず、極めて小さな株（マイクロキャップ）でも最近になって消えた。著者は、PEAD を研究する人に、それが今もあるかを確かめる方法の提案もしている。
\end{paperbox}

\begin{projbox}
LLM-BLM の \SP{} の実験で、LLM の読み取りの情報が発表の当日に株価に入りきり、翌日以降のドリフトを予測しなかったことは、Martineau（2022）の結果とよく一致する。小型株の実験 S1 は、Chordia ほか（2009）がドリフトの集中を示した流動性の低い・注目の低い銘柄に的を絞ったものである。ただし、Martineau（2022）によれば、数値の驚き（SUE）による PEAD は、極めて小さな株でも最近は消えている。S1 が問うのは、数値ではなく\textbf{文章の読み取り}が、小型株でドリフトを予測するかであり（PEAD.txt に近い問い）、先行研究から見て、予測できなくても不思議ではない。また、予測できたとしても、取引コストが高く、紙の上のドリフトが実際に取れるとは限らない（Chordia ほか 2009）。
\end{projbox}

\subsection{テキストによるドリフト}

\begin{paperbox}[Meursault・Liang・Routledge・Scanlon（2023）「PEAD.txt」\cite{peadtxt}]
決算説明会の発言録の文章から、機械学習で「文章による驚き」（SUE.txt）を作り、それで銘柄を並べると、数値の利益の驚き（SUE）によるものより大きく、長く続くドリフトが生じることを示した。数値の驚きのドリフトが近年弱まったのに対し、文章の驚きのドリフトは残っている、という。
\end{paperbox}

文章の情報は、数値の情報より処理が難しく、反映に時間がかかる、という仮説である。LLM による読み取りは、この文章の情報を安く速く処理する道具であり、その普及が、この種のドリフトを小さくしていく可能性がある（第 \ref{ch:text} 章）。

\section{アナリストの推奨の変更}

\begin{paperbox}[Womack（1996）\cite{womack1996}]
大手の証券会社のアナリストが推奨を変更すると、発表の数日に大きな反応がある。その後も同じ方向のドリフトが続き、買いへの変更の後のドリフトは小さく短い（平均 +2.4\%）が、売りへの変更の後のドリフトは大きく、6 か月にわたって続く（平均 $-9.1\%$）。
\end{paperbox}

売りの推奨の情報の反映が遅いのは、空売りの制約と、アナリストが売りの推奨を出したがらないこと（売りの推奨は情報の価値が高い）による、と解釈される。Hong・Lim・Stein（2000）\cite{hls2000} の「悪い知らせはゆっくり伝わる」とも合う。

\section{合併裁定}

\begin{paperbox}[Mitchell・Pulvino（2001）\cite{mitchellpulvino2001}]
1963〜1998 年の約 4,750 件の合併・買収で、買収の発表の後に買収される会社の株を買い（株式の交換による買収なら、買収する会社の株を売り）、成立まで持つ合併裁定の収益を調べた。合併裁定のリターンは、市場が平穏か上げているときは市場とほぼ無関係だが、市場が大きく下げるときには市場と強く連動する。この形は、指数のプット・オプションを売った戦略に似ている。この非線形なリスクを考慮しても、年 4\% 程度の超過リターンが残った。
\end{paperbox}

Baker・Savaşoglu（2002）\cite{bakersavasoglu2002} は、合併裁定の収益が、成立しないリスク（成立の失敗の確率）と、裁定に使える資金の限界（大きな案件ほど裁定者が持ちきれない）の対価であることを示した（裁定の限界）。

合併裁定の損失は、案件が不成立になったとき（独占禁止の当局の反対、資金調達の失敗、相手先の業績の悪化）に集中する。また、2008 年のような危機では、多くの案件が同時に不確かになり、裁定者が資金の制約で一斉に手じまうので、スプレッドが急に広がる（Mitchell・Pulvino 2012 \cite{mitchellpulvino2012}、第 \ref{ch:rv} 章）。

\section{指数の入れ替え}

Harris・Gurel（1986）\cite{harrisgurel1986} と Shleifer（1986）\cite{shleifer1986} は、\SP{} への採用の発表で株価が数\% 上がることを示した。Harris・Gurel は、その上昇の多くが数週間で戻ること（一時的な売買の圧力）を、Shleifer は上昇が残ること（株の需要の曲線が右下がりであること）を強調した。指数のファンドが採用の日に機械的に買うので、その前に買って指数のファンドに売る戦略が成り立つ。

Greenwood・Sammon \cite{greenwoodsammon} は、この「指数の効果」が大きく縮んだことを示した。\SP{} への採用の平均の超過リターンは、1980 年代に 3.4\%、1990 年代に 7.6\% だったが、最近の 10 年は 0.8\% にすぎない。除外も同じで、1980〜90 年代の大きな負の超過リターンが、2010〜2020 年には $-0.6\%$ になった。指数に連動する資産の割合は大きく増えたのに、効果は小さくなった（Deep Research の答え 1 の値は、NBER の要旨と一致した）。

\begin{figure}[H]
\centering
\begin{tikzpicture}[x=1mm,y=4mm]
\foreach \y in {0,2,4,6,8} {\draw[ink2!25] (0,\y)--(80,\y); \node[lab,left] at (0,\y) {\y\%};}
\fill[c1!60] (8,0) rectangle (24,3.4);  \node[lab,above] at (16,3.4) {3.4\%};
\fill[c1] (32,0) rectangle (48,7.6);    \node[lab,above] at (40,7.6) {7.6\%};
\fill[c2!75] (56,0) rectangle (72,0.8); \node[lab,above] at (64,0.8) {0.8\%};
\node[lab,below] at (16,-0.1) {1980 年代};
\node[lab,below] at (40,-0.1) {1990 年代};
\node[lab,below] at (64,-0.1) {最近の 10 年};
\end{tikzpicture}
\caption{\SP{} への採用の平均の超過リターン（Greenwood・Sammon の NBER の要旨の値から描いた）。}
\label{fig:indexeffect}
\end{figure}

\section{ディストレストとアクティビスト}

\begin{itemize}
\item \textbf{ディストレスト}：Jiang・Li・Wang（2012）\cite{jlw2012} は、米国の連邦破産法 11 条の手続きにヘッジファンドが債権者として加わると、再建が成立する確率が上がり、劣後の債権者の回収が増え、経営者の交代が増えることを示した。ヘッジファンドは、破綻した会社の債権を安く買い、手続きに影響力を持つことで収益を上げる。
\item \textbf{アクティビスト}：Brav・Jiang・Partnoy・Thomas（2008）\cite{bjpt2008} は、ヘッジファンドが 5\% 以上の株を持ったと届け出る（13D の提出）と、株価がその前後で大きな超過リターン（約 7\%）を示し、その後に戻らないこと、対象の会社の利益率や配当が改善することを示した。
\end{itemize}

どちらも、大きな資金と、法律・企業の調査の専門性が必要で、規則による再現は難しい。

\section{個人や研究で再現できるか}

\begin{itemize}
\item PEAD や推奨の変更の検証には、決算の日付と予想（アナリストの予想のデータベース）、あるいは決算の数値の時系列が要る。前者は有料、後者は SEC の XBRL などで無料に近い形で作れる。
\item 近年の大型株では PEAD がほぼ消えているので、検証するなら、注目・流動性の低い銘柄や、文章の情報（PEAD.txt）に的を絞ることになる。その場合、取引コストを厳しく見積もる必要がある。
\item 合併裁定は、発表された案件の一覧と価格があれば検証できる。不成立の案件の扱い（生存者バイアス）に注意する。
\end{itemize}
